\documentclass[10pt,a4paper]{amsart}
\usepackage[margin=19mm]{geometry}
\usepackage{amsmath,amssymb,mathtools}
\usepackage[scr=rsfs,cal=euler]{mathalfa}
\usepackage{xfrac}
\usepackage[unicode,hidelinks,bookmarks=false]{hyperref}
\newcommand{\ie}{i.e.\ }
\newcommand{\cf}{cf.\ }
\newcommand{\resp}{respectively\ }
\newcommand{\Subsection}{\subsection}
\providecommand{\nobreakdash}{\nobreak\hbox{-}}
\setlength{\emergencystretch}{2em}
\multlinegap0pt
\usepackage{float}
\usepackage{amssymb}
\usepackage{mathtools}
\usepackage[normalem]{ulem}
\newtheorem{assumption}{Assumption}
\newtheorem{lemma}{Lemma}
\newtheorem{proposition}{Proposition}
\title{Sobolev Regularized Score Difference Estimation in Diffusion Models}
\author{Chenghan Xie, Jose Blanchet, Renyuan Xu}
\date{Source: arXiv:2608.18237v2 (CC BY 4.0)}
\begin{document}
\maketitle

\noindent\emph{Source and licence.} Sobolev Regularized Score Difference Estimation in Diffusion Models. Version arXiv:2608.18237v2 (CC BY 4.0), distributed by arXiv under the Creative Commons Attribution 4.0 International licence, http://creativecommons.org/licenses/by/4.0/, as recorded in the arXiv licence metadata for this exact version and shown on the abstract page https://arxiv.org/abs/2608.18237v2. Full licence notice: \texttt{licenses/arxiv-2608.18237-CC-BY-4.0.txt}.\par\smallskip

\noindent\textit{Editorial scope.} Counted unit: the article's proposition NN (source0.tex lines 1279--1321). The article's complete original proof of it is delivered with it (source0.tex lines 1323--1631, one \texttt{proof} environment of 309 lines). No mathematics is written, repaired or re-derived: the statement and the proof are the article's own lines, in the article's own notation, and no retained line is substituted or re-flowed.  Retained uncounted as the source-local context the proof needs: lines 229--241; lines 1060--1104.  Nothing was omitted from any retained span, so the package carries no omission record.  No retained line contains a citation, so the wrapper carries no bibliography and no bibliography entry.  No retained line contains a figure environment, an image-inclusion command or drawing code, so no image is delivered and the package declares no asset dependency.  Editorial formatting only, in addition: an AMS \texttt{amsart} preamble that restates the article's own declarations and macros used by the retained lines, namely the environments \texttt{assumption}, \texttt{lemma}, \texttt{proposition} on separate counters, together with the standard packages those lines need.  The retained environments print in this document as Assumption 1, Lemma 1, Proposition 1 in order of appearance, whereas the article prints the same items with the numbering of its own full document.  The complete native source of the article as distributed in the arXiv e-print for this exact version is bundled unchanged as \texttt{sources/207977/source0.tex}.
\begin{assumption}[Density functions]\label{assm: density}
Assume  $p$ and $q$ satisfy the following conditions:
\begin{enumerate}
\item $p,q\in C^2(\overline\Omega)$, and there exist constants \(0<c<C<\infty\) such that \(c\le p(x),q(x)\le C\) for all \(x\in\overline\Omega\).
\item The following Neumann boundary condition holds:
\[
\rho(x)\,\partial_n f^\star(x)=0
\qquad \text{for all }x\in\partial\Omega,
\]
where \(n(x)\) is the outward unit normal to \(\partial\Omega\).
\item Fix $s\geq2$. We assume $f^\star \in H^s(\mu_0)$.
\end{enumerate}
\end{assumption}

\begin{lemma}[Tensor-product B-spline quasi-interpolation]
\label{lem:spline-stability}
Fix integers \(k,l\in\mathbb N\). There exist
\[
Q_{k,l}:L^1([0,1]^d)\to L^1([0,1]^d),
\qquad
C_{k,d}=C(k,d)<\infty,
\qquad
C_q=C_q(k,d)<\infty,
\]
such that
\[
Q_{k,l}g(x)
=
Q_{k,l}\!\left(
g\,\mathbf 1_{\{u\in[0,1]^d:\|u-x\|_\infty\le C_{k,d}/l\}}
\right)(x),
\qquad
g\in L^1([0,1]^d),\ x\in[0,1]^d,
\]
and
\begin{equation}
\label{eq:qi-Linfty-gradient-stability}
\|Q_{k,l}g\|_{L^\infty}
\le
C_q\|g\|_{L^\infty},
\qquad
\|\nabla Q_{k,l}g\|_{L^\infty}
\le
C_q\|\nabla g\|_{L^\infty},
\qquad
g,\partial_1g,\ldots,\partial_dg\in L^\infty([0,1]^d).
\end{equation}
If \(s\in\mathbb N\), \(0\le r\le s\), and \(k\ge s\), then, for
\(C=C(k,s,r,d)\),
\begin{equation}
\label{eq:bspline-Hr-approx-general}
\|Q_{k,l}g-g\|_{H^r}
\le
C l^{-(s-r)}
\|g\|_{H^s},
\qquad
g\in H^s([0,1]^d).
\end{equation}
\end{lemma}

\begin{proposition}[Approximation by the clipped and gradient-bounded sieve]
\label{prop: approximate of clipped NN}
Let \(\Omega\subset\mathbb R^d\) be a bounded domain with sufficiently smooth
boundary, and let \(\mu\) be the spatial mixture measure corresponding to a pair
\((p,q)\) satisfying Assumption~\ref{assm: density}. Let \(s\in\mathbb N\)
with \(1\le s\le4\). Suppose
$
f^\star\in H^s(\mu)
$
and
\[
\max\left\{
\|f^\star\|_{L^\infty(\mu)},
\|\nabla f^\star\|_{L^\infty(\mu)}
\right\}
\le M^\star .
\]
Then there exists \(C_\Omega\ge1\), depending only on \(\Omega\), the density
bounds in Assumption~\ref{assm: density}, and the fixed dimension \(d\), such
that the following holds. Choose \(M>C_\Omega M^\star\), and choose \(T_M\)
so that \(T_M(u)=u\) for all \(|u|\le C_\Omega M^\star\). For
\[
L=\mathcal O(1),
\qquad
W=\mathcal O(N),
\qquad
S=\mathcal O(N),
\qquad
B=\mathcal O(N),
\]
there exists
\[
\bar f_N\in\mathcal F_M(L,W,S,B)
\]
such that
\begin{equation}
\label{eq:clipped-approx-H1-final}
\|\bar f_N-f^\star\|_{H^1(\mu)}^2
\lesssim
N^{-\frac{2(s-1)}{d}}
\|f^\star\|_{H^s(\mu)}^2 .
\end{equation}
\end{proposition}

\begin{proof}
Under Assumption~\ref{assm: density}, the weighted Sobolev norm induced by
\(\mu\) is equivalent to the usual Sobolev norm on \(\Omega\). Hence it
suffices to prove the approximation result in the usual \(H^1(\Omega)\) norm.

Let
\[
\mathcal C_\Omega
=
\prod_{j=1}^d[\alpha_j,\alpha_j+\rho_\Omega]
\]
be a closed axis-aligned cube containing \(\overline\Omega\), with
\(\rho_\Omega>0\). Since \(\Omega\) has sufficiently smooth boundary, choose a
bounded linear extension operator
\[
\mathcal E:H^s(\Omega)\to H^s(\mathcal C_\Omega)
\]
such that, for a constant \(C_{\mathcal E}\ge1\),
\[
\|\mathcal E f\|_{H^s(\mathcal C_\Omega)}
\le
C_{\mathcal E}\|f\|_{H^s(\Omega)}
\]
and
\[
\max\left\{
\|\mathcal E f\|_{L^\infty(\mathcal C_\Omega)},
\|\nabla(\mathcal E f)\|_{L^\infty(\mathcal C_\Omega)}
\right\}
\le
C_{\mathcal E}
\max\left\{
\|f\|_{L^\infty(\Omega)},
\|\nabla f\|_{L^\infty(\Omega)}
\right\}.
\]
Define the affine map
\[
a_\Omega(x)
:=
\left(
\frac{x_1-\alpha_1}{\rho_\Omega},
\ldots,
\frac{x_d-\alpha_d}{\rho_\Omega}
\right),
\qquad x\in \mathcal C_\Omega .
\]
Then \(a_\Omega\) is a bijection from \(\mathcal C_\Omega\) to \([0,1]^d\)
with linear part \(R_\Omega=\rho_\Omega^{-1}I_d\), and hence \(R_\Omega\) is
invertible. Moreover,
\[
a_\Omega^{-1}(z)
=
(\alpha_1+\rho_\Omega z_1,\ldots,\alpha_d+\rho_\Omega z_d),
\qquad z\in[0,1]^d .
\]
Let \(C_q\) be the stability constant in
Lemma~\ref{lem:spline-stability} for cubic splines, and set
\[
C_\Omega
:=
\max\left\{
1,\,
C_qC_{\mathcal E},\,
\|R_\Omega\|_{\mathrm{op}}C_q
\|R_\Omega^{-1}\|_{\mathrm{op}}C_{\mathcal E}
\right\}.
\]
All constants below may depend on \(\Omega\), the density bounds in
Assumption~\ref{assm: density}, the extension operator, the affine map
\(a_\Omega\), and the fixed dimension \(d\), but not on \(N\).

Let
\[
l:=\left\lceil N^{1/d}\right\rceil .
\]
Define the spline approximant on \(\Omega\) by
\begin{equation}\label{eq:def of F_N}
    F_N(x)
:=
\left(
Q_{4,l}\bigl((\mathcal E f^\star)\circ a_\Omega^{-1}\bigr)
\right)(a_\Omega(x)),
\qquad x\in\Omega .
\end{equation}
Here \((\mathcal E f^\star)\circ a_\Omega^{-1}\) is a function on
\([0,1]^d\), so \(Q_{4,l}\) is applied on the unit cube.

We first prove the approximation bound. Since \(a_\Omega\) is fixed, Sobolev
norms before and after the affine change of variables are equivalent.
Therefore
\[
\|(\mathcal E f^\star)\circ a_\Omega^{-1}\|_{H^s([0,1]^d)}
\lesssim
\|\mathcal E f^\star\|_{H^s(\mathcal C_\Omega)}
\le
C_{\mathcal E}\|f^\star\|_{H^s(\Omega)}
\lesssim
\|f^\star\|_{H^s(\mu)}.
\]
By Lemma~\ref{lem:spline-stability}, with \(k=4\) and \(r=1\),
\[
\left\|
Q_{4,l}\bigl((\mathcal E f^\star)\circ a_\Omega^{-1}\bigr)
-
(\mathcal E f^\star)\circ a_\Omega^{-1}
\right\|_{H^1([0,1]^d)}
\lesssim
l^{-(s-1)}
\|(\mathcal E f^\star)\circ a_\Omega^{-1}\|_{H^s([0,1]^d)}.
\]
Since \(a_\Omega\) is fixed and \(\mathcal E f^\star=f^\star\) on
\(\Omega\), by the definition of $F_N$ \eqref{eq:def of F_N}, the same change of variables gives
\[
\|F_N-f^\star\|_{H^1(\Omega)}
\lesssim
\left\|
Q_{4,l}\bigl((\mathcal E f^\star)\circ a_\Omega^{-1}\bigr)
-
(\mathcal E f^\star)\circ a_\Omega^{-1}
\right\|_{H^1([0,1]^d)}.
\]
Combining the previous estimates,
\[
\|F_N-f^\star\|_{H^1(\Omega)}
\lesssim
l^{-(s-1)}
\|f^\star\|_{H^s(\mu)}.
\]
Since \(l\asymp N^{1/d}\),
\[
\|F_N-f^\star\|_{H^1(\Omega)}
\lesssim
N^{-\frac{s-1}{d}}
\|f^\star\|_{H^s(\mu)}.
\]
By the norm equivalence between \(H^1(\Omega)\) and \(H^1(\mu)\), this
also implies
\[
\|F_N-f^\star\|_{H^1(\mu)}
\lesssim
N^{-\frac{s-1}{d}}
\|f^\star\|_{H^s(\mu)}.
\]

We next verify the uniform envelope. By the boundedness of the extension
operator,
\[
\|(\mathcal E f^\star)\circ a_\Omega^{-1}\|_{L^\infty([0,1]^d)}
=
\|\mathcal E f^\star\|_{L^\infty(\mathcal C_\Omega)}
\le
C_{\mathcal E}M^\star .
\]
By the \(L^\infty\)-stability of \(Q_{4,l}\) in Lemma \ref{lem:spline-stability},
\[
\|F_N\|_{L^\infty(\Omega)}
\le
\left\|
Q_{4,l}\bigl((\mathcal E f^\star)\circ a_\Omega^{-1}\bigr)
\right\|_{L^\infty([0,1]^d)}
\le
C_qC_{\mathcal E}M^\star
\le
C_\Omega M^\star
<
M.
\]

For the gradient, the chain rule gives
\[
\nabla\bigl((\mathcal E f^\star)\circ a_\Omega^{-1}\bigr)(z)
=
R_\Omega^{-\top}
\nabla(\mathcal E f^\star)(a_\Omega^{-1}(z)).
\]
Hence
\[
\left\|
\nabla\bigl((\mathcal E f^\star)\circ a_\Omega^{-1}\bigr)
\right\|_{L^\infty([0,1]^d)}
\le
\|R_\Omega^{-1}\|_{\mathrm{op}}C_{\mathcal E}M^\star .
\]
Using the gradient stability of \(Q_{4,l}\) in Lemma \ref{lem:spline-stability},
\[
\left\|
\nabla Q_{4,l}\bigl((\mathcal E f^\star)\circ a_\Omega^{-1}\bigr)
\right\|_{L^\infty([0,1]^d)}
\le
C_q
\left\|
\nabla\bigl((\mathcal E f^\star)\circ a_\Omega^{-1}\bigr)
\right\|_{L^\infty([0,1]^d)}.
\]
Since
\[
F_N(x)
=
\left(
Q_{4,l}\bigl((\mathcal E f^\star)\circ a_\Omega^{-1}\bigr)
\right)(a_\Omega(x)),
\]
another application of the chain rule gives
\[
\|\nabla F_N\|_{L^\infty(\Omega)}
\le
\|R_\Omega\|_{\mathrm{op}}
\left\|
\nabla Q_{4,l}\bigl((\mathcal E f^\star)\circ a_\Omega^{-1}\bigr)
\right\|_{L^\infty([0,1]^d)}.
\]
Combining the last three inequalities,
\[
\|\nabla F_N\|_{L^\infty(\Omega)}
\le
\|R_\Omega\|_{\mathrm{op}}
C_q
\|R_\Omega^{-1}\|_{\mathrm{op}}
C_{\mathcal E}M^\star
\le
C_\Omega M^\star
<
M.
\]
Thus the definition of \(C_\Omega\) ensures that both the function value and
the gradient of \(F_N\) are controlled before clipping.

It remains to realize \(F_N\) as a sparse ReLU\(^3\) neural network. The
quasi-interpolant
\[
Q_{4,l}\bigl((\mathcal E f^\star)\circ a_\Omega^{-1}\bigr)
\]
is a linear combination of tensor-product cubic B-splines on \([0,1]^d\).
For interior indices,
\[
N_{l,i}^{(4)}(z)
=
\frac{1}{3!}
\sum_{j=0}^{4}
(-1)^j
\binom{4}{j}
\left(lz-(i+j)\right)_+^3 .
\]
The boundary splines have analogous formulas determined by the repeated
endpoint knots. Hence every univariate cubic B-spline can be represented
exactly by a ReLU\(^3\) subnetwork. Since the dimension \(d\) is fixed,
each tensor-product cubic B-spline can be represented by a fixed-depth
ReLU\(^3\) subnetwork. The number of tensor-product basis functions is
\[
|I_{l,4}|^d=\mathcal O(l^d)=\mathcal O(N).
\]
Composing with the fixed affine map \(a_\Omega\) only modifies the first layer.
Therefore \(F_N\) can be represented by a sparse ReLU\(^3\) network with
\[
L=\mathcal O(1),
\qquad
W=\mathcal O(N),
\qquad
S=\mathcal O(N),
\qquad
B=\mathcal O(N).
\]
The constants hidden in the \(\mathcal O(\cdot)\) notation may depend on
\(\Omega\), \(d\), the spline order, \(M^\star\), and fixed problem
parameters, but not on \(N\).

Finally define
\[
\bar f_N:=T_M\circ F_N .
\]
Since
\[
\|F_N\|_{L^\infty(\Omega)}\le C_\Omega M^\star
\]
and \(T_M(u)=u\) for all \(|u|\le C_\Omega M^\star\), the clipping is inactive
on \(\Omega\). Therefore
\[
\bar f_N=F_N
\qquad\text{on }\Omega .
\]
Consequently,
\[
\|\bar f_N\|_{L^\infty(\mu)}\le M,
\qquad
\|\nabla \bar f_N\|_{L^\infty(\mu)}\le M,
\]
and hence
\[
\bar f_N\in\mathcal F_M(L,W,S,B).
\]
Moreover,
\[
\|\bar f_N-f^\star\|_{H^1(\mu)}
=
\|F_N-f^\star\|_{H^1(\mu)}
\lesssim
N^{-\frac{s-1}{d}}
\|f^\star\|_{H^s(\mu)}.
\]
Squaring both sides gives
\[
\|\bar f_N-f^\star\|_{H^1(\mu)}^2
\lesssim
N^{-\frac{2(s-1)}{d}}
\|f^\star\|_{H^s(\mu)}^2 .
\]
This proves \eqref{eq:clipped-approx-H1-final}.
\end{proof}

\end{document}
